\section*{NeurIPS Paper Checklist}

\begin{enumerate}

\item {\bf Claims}
    \item[] Question: Do the main claims made in the abstract and introduction accurately reflect the paper's contributions and scope?
    \item[] Answer: \answerYes{}
    \item[] Justification: The abstract enumerates the contributions
          (the \bench dataset, the three-task evaluation protocol,
          the 16-model evaluation, and the five-axis analysis) and
          the scope (single-image anomaly detection in dashcam
          driving scenes), all of which are explicitly delivered in
          Sections~\ref{sec:dataset}--\ref{sec:analysis}.

\item {\bf Limitations}
    \item[] Question: Does the paper discuss the limitations of the work performed by the authors?
    \item[] Answer: \answerYes{}
    \item[] Justification: Section~\ref{sec:discussion} contains a
          dedicated ``Limitations'' paragraph covering: (i)~auto-generated
          GPT-4o annotations and the cross-LLM validation we performed
          to mitigate this; (ii)~single-image scope (no temporal
          anomalies); (iii)~focus on open-source VLMs only;
          (iv)~zero-shot prompting versus fine-tuning.

\item {\bf Theory assumptions and proofs}
    \item[] Question: For each theoretical result, does the paper provide the full set of assumptions and a complete (and correct) proof?
    \item[] Answer: \answerNA{}
    \item[] Justification: The paper is an empirical
          benchmark and analysis paper; it contains no theoretical
          theorems or proofs.

    \item {\bf Experimental result reproducibility}
    \item[] Question: Does the paper fully disclose all the information needed to reproduce the main experimental results of the paper to the extent that it affects the main claims and/or conclusions of the paper (regardless of whether the code and data are provided or not)?
    \item[] Answer: \answerYes{}
    \item[] Justification: All prompts are reproduced in
          Appendix~\ref{app:prompts}; all model identifiers, hyperparameters,
          decoding settings (temperature, max new tokens), prompt-caching
          choices, and resource budgets are documented. Code and data are
          released with the paper (Appendices~\ref{app:reproducibility}
          and~\ref{app:license}).

\item {\bf Open access to data and code}
    \item[] Question: Does the paper provide open access to the data and code, with sufficient instructions to faithfully reproduce the main experimental results, as described in supplemental material?
    \item[] Answer: \answerYes{}
    \item[] Justification: \bench is released under
          CC-BY 4.0 with a datasheet (Appendix~\ref{app:license}). The
          full evaluation and analysis pipeline is open-sourced
          (Appendix~\ref{app:reproducibility}) with exact reproduction
          commands.

\item {\bf Experimental setting/details}
    \item[] Question: Does the paper specify all the training and test details (e.g., data splits, hyperparameters, how they were chosen, type of optimizer) necessary to understand the results?
    \item[] Answer: \answerYes{}
    \item[] Justification: We do not train any new VLMs;
          all VLMs are evaluated zero-shot with prompts in
          Appendix~\ref{app:prompts}. The autoencoder baseline is
          described in Sec.~\ref{sec:models} with hyperparameters in
          the released code. CLIP/SigLIP prompts are listed in
          Appendix~\ref{app:prompts}.

\item {\bf Experiment statistical significance}
    \item[] Question: Does the paper report error bars suitably and correctly defined or other appropriate information about the statistical significance of the experiments?
    \item[] Answer: \answerYes{}
    \item[] Justification: The class-difficulty bar
          chart (Fig.~\ref{fig:hardest_classes}) and the consistency
          analysis (Sec.~\ref{sec:consistency},
          Table~\ref{tab:consistency_multiclass}) report error bars
          (standard deviation across models / runs). Headline
          benchmark numbers in Tables~\ref{tab:binary_results} and
          \ref{tab:multiclass_results} are deterministic at $T=0$.

\item {\bf Experiments compute resources}
    \item[] Question: For each experiment, does the paper provide sufficient information on the computer resources (type of compute workers, memory, time of execution) needed to reproduce the experiments?
    \item[] Answer: \answerYes{}
    \item[] Justification: Appendix~\ref{app:compute}
          itemises GPU-hours per experiment phase (binary,
          multiclass, description, consistency, hidden-state
          extraction) and total ($\sim$845 GPU-hours on RTX 5090s).

\item {\bf Code of ethics}
    \item[] Question: Does the research conducted in the paper conform, in every respect, with the NeurIPS Code of Ethics?
    \item[] Answer: \answerYes{}
    \item[] Justification: We have read the NeurIPS Code
          of Ethics and confirm full conformance. All data was
          collected from publicly licensed sources; faces and license
          plates were screened/removed; no personally identifiable
          information appears in the dataset.

\item {\bf Broader impacts}
    \item[] Question: Does the paper discuss both potential positive societal impacts and negative societal impacts of the work performed?
    \item[] Answer: \answerYes{}
    \item[] Justification: Section~\ref{sec:discussion}
          contains an explicit broader-impacts subsection.
          Positive: the benchmark accelerates safer VLM-based
          perception research. Negative: misinterpretation as an
          endorsement to deploy current VLMs in safety-critical AD;
          we explicitly counsel against this use in the dataset
          documentation.

\item {\bf Safeguards}
    \item[] Question: Does the paper describe safeguards that have been put in place for responsible release of data or models that have a high risk for misuse (e.g., pre-trained language models, image generators, or scraped datasets)?
    \item[] Answer: \answerYes{}
    \item[] Justification: We release no new model weights,
          only image data and annotations. The dataset documentation
          (datasheet) explicitly warns against safety-critical
          deployment based on the released models' performance. All
          potentially-identifiable visual content (faces, plates) was
          screened before release.

\item {\bf Licenses for existing assets}
    \item[] Question: Are the creators or original owners of assets (e.g., code, data, models), used in the paper, properly credited and are the license and terms of use explicitly mentioned and properly respected?
    \item[] Answer: \answerYes{}
    \item[] Justification: All evaluated VLMs and CLIP/SigLIP
          variants are cited with their canonical references. Image
          source datasets and their licenses are listed in
          Appendix~\ref{app:license}; all sources are licensed in a
          manner compatible with CC-BY redistribution.

\item {\bf New assets}
    \item[] Question: Are new assets introduced in the paper well documented and is the documentation provided alongside the assets?
    \item[] Answer: \answerYes{}
    \item[] Justification: \bench is released with a
          datasheet~\citep{datasheets} covering motivation,
          composition, collection process, recommended uses,
          maintenance plan, and known limitations.

\item {\bf Crowdsourcing and research with human subjects}
    \item[] Question: For crowdsourcing experiments and research with human subjects, does the paper include the full text of instructions given to participants and screenshots, if applicable, as well as details about compensation (if any)?
    \item[] Answer: \answerNA{}
    \item[] Justification: No crowdsourced annotation was
          used. Annotations were generated by GPT-4o, validated by
          Claude Sonnet 4.6, and manually adjudicated by the authors.

\item {\bf Institutional review board (IRB) approvals or equivalent for research with human subjects}
    \item[] Question: Does the paper describe potential risks incurred by study participants, whether such risks were disclosed to the subjects, and whether Institutional Review Board (IRB) approvals (or an equivalent approval/review based on the requirements of your country or institution) were obtained?
    \item[] Answer: \answerNA{}
    \item[] Justification: The work does not involve human
          subjects research.

\item {\bf Declaration of LLM usage}
    \item[] Question: Does the paper describe the usage of LLMs if it is an important, original, or non-standard component of the core methods in this research?
    \item[] Answer: \answerYes{}
    \item[] Justification: Appendix~\ref{app:llm_usage}
          declares the use of GPT-4o for initial annotation
          generation and Claude Sonnet 4.6 for cross-validation;
          both are described in detail in
          Sec.~\ref{sec:dataset} and
          Appendix~\ref{app:validation}.

\end{enumerate}
